\section{Tổ chức nội dung: gói giáo án và đồ thị tri thức}\label{sec:content-design}
Một chương được biểu diễn bằng đồ thị có hướng $G=(V,E)$. Mỗi đỉnh là một bài học, mỗi cạnh $(u,v)$ mang một trong hai kiểu: tiên quyết (học $u$ trước $v$) hoặc vận dụng ($v$ vận dụng $u$). Hai chương Toán có $19$ bài và $28$ cạnh: chương Số thập phân có $11$ bài và $17$ cạnh, chương Diện tích có $8$ bài và $11$ cạnh; Unit~11 Tiếng Anh có $6$ bài. Mỗi đỉnh đi kèm một gói nội dung gồm các cách giảng lý thuyết, phương pháp làm bài, ví dụ mẫu, ghi chú lỗi sai thường gặp, ngân hàng câu hỏi và các ý chính của bài để học sinh kể lại. Một câu hỏi có đề bài và đặc tả đáp án (số, trắc nghiệm, so sánh, sắp xếp, đúng sai hoặc mở), mức nhận thức (A đến E và chẩn đoán), độ khó từ $0$ đến $1$, mục đích (kiểm tra hiểu, luyện tập, kiểm tra cuối bài, ôn tập) và thang gợi ý riêng. Mức nhận thức và độ khó quyết định chiều của mô hình người học mà câu trả lời cung cấp bằng chứng (mục~\ref{sec:learner-design}).

Gói giáo án được dựng từ $18$ tài liệu thật của giáo viên ($77$ trang: giáo án, sách giáo khoa, sách bài tập, phiếu bài tập), không soạn tay. Ba luồng chạy song song trên văn bản đã đọc: dựng lược đồ bài học, cắt từng câu hỏi và đáp án, và trích nội dung từng bài. Mọi đáp án trích được đều được kiểm lại bằng cách tính lại từ đề; đáp án mâu thuẫn với đề bị loại và được liệt kê trong báo cáo thay vì được dạy. Ngân hàng của chương Số thập phân có $234$ câu ($46$ câu kiểm tra), của chương Diện tích có $155$ câu ($33$ câu kiểm tra). Phần tài liệu còn thiếu được mô hình soạn nháp, chỉ dùng sau khi qua kiểm chứng và được giáo viên duyệt, và nội dung của giáo viên luôn thắng nội dung do mô hình đề xuất. Quyết định duyệt gắn với phiên bản xuất bản của kế hoạch, nên mỗi buổi học dùng đúng phiên bản chương lúc bắt đầu và một buổi cũ phát lại luôn cho cùng kết quả.

Cùng nội dung được cắt thành các đoạn truy xuất. Chỉ mục xuất bản của hai chương Toán có $221$ đoạn: $83$ đoạn trích từ tài liệu nguồn, $38$ đoạn lý thuyết, $27$ ví dụ mẫu, $19$ đoạn phương pháp và $54$ ghi chú lỗi quan niệm. Mỗi đoạn mang mã bài của nó trong đồ thị, nên truy xuất biết câu hỏi được đặt khi học sinh đang ở bài nào và biết khoảng cách trên đồ thị đến bài chứa đoạn (mục~\ref{sec:retrieval-design}).
